\scanpage{14}
\mspage{6}
\margin{c'est le nb de pts à l'infini de la composante $X_\alpha$}
$\operatorname{card}I_\alpha$ (où $I_\alpha$ est la fibre de $I$ en $\alpha\in S$)
\[
(16)\qquad\nu_\alpha=\operatorname{card}I_\alpha,
\]
et on appelle \uline{poids total} d'un sommet l'entier
\[
(17)\qquad\widehat\nu_\alpha=\nu_\alpha+
\underset{\substack{\Vert\ \mathrm{déf}\\
\text{ordre du sommet }\alpha,\\
\text{i.e. }\operatorname{card}\widetilde A_\alpha,\text{ où }\widetilde A_\alpha\text{ est la fibre}\\
\text{de }\widetilde A\text{ au dessus de }\alpha,\text{ i.e. le}\\
\text{nb d'arcs sortant de }\alpha}}{\omega_\alpha}
\]
c'est donc le nb de ``pts marqués'' (provenant de pts doubles \unsure{et} de pts à l'infini de $X$) sur la composante $Y_\alpha$ de $Y=\widetilde{\widehat X}$.

3. \uline{Courbes ``stables''.} \uline{et} \uline{MD-graphes.}

Une courbe \unsure{standard|standart} \inserted{(sur $k$ alg. clos)} est ``stable'', si elle satisfait à l'une des conditions équivalentes suivantes

a) $\operatorname{Aut}X$ est fini

b) Pour tout $\alpha$, $(Y_\alpha,I_\alpha\cup\widetilde A_\alpha)$ est anabélien
\[
\text{i.e.}\quad 2g_\alpha+\widehat\nu_\alpha\geq3\qquad
\text{i.e.}\quad 2g_\alpha-2+\widehat\nu_\alpha\geq1,\quad\text{i.e.}
\]
\mspage{13}
